% GENERATED FILE. Edit canonical lessons, then run npm run generate:books.
% canonical-source-hash: fnv1a64:f5aa1b7d7f7c6254
% canonical-lessons: HI-C292-daku, HI-C292-lutna, HI-C292-dakaiti, HI-C292-banduq, HI-C292-badhna

\chapter{Robbers on the Road}
\label{ch:hi-robbers-on-the-road}
\hlchaptermodality{292}

\begin{chapteropening}
\textbf{By the end of this chapter:} \emph{I can tell what robbers did and say what will happen if they come back.}

Complete the last lesson of chapter 292: I can tell what robbers did and say what will happen if they come back.
\end{chapteropening}

\section[\texorpdfstring{\hi{डाकू} (ḍākū)}{डाकू (ḍākū)}]{\hi{डाकू} (ḍākū) — a bandit}

\label{lesson:HI-C292-daku}



\begin{quote}
Before the new one: say the Hindi for suspicion, then the Hindi for to stare.
\end{quote}

\subsection*{You'll want to know: \hi{डाकू}}
\textbf{\hi{डाकू}} — \emph{ḍākū} — \textquotedblleft{}a bandit\textquotedblright{}.

\begin{quote}
\hi{रात} \hi{को} \hi{दरवाज़ा} \hi{बंद} \hi{रखो}, \hi{नहीं} \hi{तो} \hi{डाकू} \hi{आ} \hi{जाएँगे।}
\end{quote}

\begin{quote}
\emph{rāt ko darvāzā band rakho, nahīṃ to ḍākū ā jāeṅge.}
\end{quote}

\begin{quote}
Keep the door shut at night, or the bandits will get in.
\end{quote}

Masculine.

\begin{cousinweb}
comes from \emph{*ḍākka}, \textquotedblleft{}robbery\textquotedblright{}, a form CDIAL reconstructs from its descendants (CDIAL 5543). A \textbf{\hi{डाकू}} belongs to a gang that robs on the roads and in villages, not a single \textbf{\hi{चोर}} who steals in secret.
\end{cousinweb}

\begin{grammarlens}[title={What you've built}]
One more word for telling what happened.
\end{grammarlens}

\subsection*{Your turn}
\begin{itemize}
\raggedright
  \item \emph{Say it:} \emph{ḍākū}
  \item \emph{Say it:} \emph{ḍākū}, once more
  \item \emph{Say it:} \emph{rāt ko darvāzā band rakho, nahīṃ to ḍākū ā jāeṅge.} — Keep the door shut at night, or the bandits will get in.
  \item \emph{Recall:} say the Hindi for suspicion, then the Hindi for to stare, then say \emph{ḍākū} again
\end{itemize}

\begin{tcolorbox}[breakable,colback=teal!4,colframe=teal!35!black,title={Before you move on}]
What does \hi{डाकू} mean? (\textquotedblleft{}A bandit\textquotedblright{}.) Say it once more.
\end{tcolorbox}
\section[\texorpdfstring{\hi{लूटना} (lūṭnā)}{लूटना (lūṭnā)}]{\hi{लूटना} (lūṭnā) — to rob}

\label{lesson:HI-C292-lutna}



\begin{quote}
Before the new one: say the Hindi for to stare, then the Hindi for a bandit.
\end{quote}

\subsection*{You'll want to know: \hi{लूटना}}
\textbf{\hi{लूटना}} — \emph{lūṭnā} — \textquotedblleft{}to rob\textquotedblright{}.

\begin{quote}
\hi{डाकुओं} \hi{ने} \hi{पूरा} \hi{गाँव} \hi{लूटा।}
\end{quote}

\begin{quote}
\emph{ḍākuoṁ ne pūrā gā̃v lūṭā.}
\end{quote}

\begin{quote}
The bandits robbed the whole village.
\end{quote}

Past, with \hi{ने} (he or she …): \hi{उसने} \hi{दुकान} \hi{लूटी}, \emph{usne dukān lūṭī}.

\begin{cousinweb}
\textbf{\hi{लूटना}} (\emph{lūṭnā}) is inherited from \emph{*luṭṭati}, \textquotedblleft{}plunders\textquotedblright{}, a form CDIAL reconstructs from its descendants (CDIAL 11078). English \emph{loot} was taken from this family of words.
\end{cousinweb}

\begin{grammarlens}[title={What you've built}]
One more word for telling what happened.
\end{grammarlens}

\subsection*{Your turn}
\begin{itemize}
\raggedright
  \item \emph{Say it:} \emph{lūṭnā}
  \item \emph{Say it:} \emph{lūṭnā}, once more
  \item \emph{Say it:} \emph{ḍākuoṁ ne pūrā gā̃v lūṭā.} — The bandits robbed the whole village.
  \item \emph{Recall:} say the Hindi for to stare, then the Hindi for a bandit, then say \emph{lūṭnā} again
\end{itemize}

\begin{tcolorbox}[breakable,colback=teal!4,colframe=teal!35!black,title={Before you move on}]
What does \hi{लूटना} mean? (\textquotedblleft{}To rob\textquotedblright{}.) Say it once more.
\end{tcolorbox}
\section[\texorpdfstring{\hi{डकैती} (ḍakaitī)}{डकैती (ḍakaitī)}]{\hi{डकैती} (ḍakaitī) — a robbery (by a gang)}

\label{lesson:HI-C292-dakaiti}



\begin{quote}
Before the new one: say the Hindi for a bandit, then the Hindi for to rob.
\end{quote}

\subsection*{You'll want to know: \hi{डकैती}}
\textbf{\hi{डकैती}} — \emph{ḍakaitī} — \textquotedblleft{}a robbery (by a gang)\textquotedblright{}.

\begin{quote}
\hi{अगर} \hi{फिर} \hi{डकैती} \hi{हुई} \hi{तो} \hi{गाँव} \hi{के} \hi{लोग} \hi{पुलिस} \hi{को} \hi{बुलाएँगे।}
\end{quote}

\begin{quote}
\emph{agar phir ḍakaitī huī to gā̃v ke log pulis ko bulāeṅge.}
\end{quote}

\begin{quote}
If there is another robbery, the villagers will call the police.
\end{quote}

Feminine.

\begin{cousinweb}
\textbf{\hi{डकैती}} (\emph{ḍakaitī}) is \emph{ḍakait}, a member of a robber gang, with the noun ending \emph{-ī}; CDIAL 5543 lists \emph{ḍakait} beside \textbf{\hi{डाकू}} under \emph{*ḍākka}, \textquotedblleft{}robbery\textquotedblright{}. English \emph{dacoit} and \emph{dacoity} were taken from these words.
\end{cousinweb}

\begin{grammarlens}[title={What you've built}]
One more word for telling what happened.
\end{grammarlens}

\subsection*{Your turn}
\begin{itemize}
\raggedright
  \item \emph{Say it:} \emph{ḍakaitī}
  \item \emph{Say it:} \emph{ḍakaitī}, once more
  \item \emph{Say it:} \emph{agar phir ḍakaitī huī to gā̃v ke log pulis ko bulāeṅge.} — If there is another robbery, the villagers will call the police.
  \item \emph{Recall:} say the Hindi for a bandit, then the Hindi for to rob, then say \emph{ḍakaitī} again
\end{itemize}

\begin{tcolorbox}[breakable,colback=teal!4,colframe=teal!35!black,title={Before you move on}]
What does \hi{डकैती} mean? (\textquotedblleft{}A robbery (by a gang)\textquotedblright{}.) Say it once more.
\end{tcolorbox}
\section[\texorpdfstring{\hi{बंदूक़} (bandūq)}{बंदूक़ (bandūq)}]{\hi{बंदूक़} (bandūq) — a gun}

\label{lesson:HI-C292-banduq}



\begin{quote}
Before the new one: say the Hindi for to rob, then the Hindi for a robbery (by a gang).
\end{quote}

\subsection*{You'll want to know: \hi{बंदूक़}}
\textbf{\hi{बंदूक़}} — \emph{bandūq} — \textquotedblleft{}a gun\textquotedblright{}.

\begin{quote}
\hi{एक} \hi{डाकू} \hi{के} \hi{हाथ} \hi{में} \hi{बंदूक़} \hi{थी।}
\end{quote}

\begin{quote}
\emph{ek ḍākū ke hāth meṁ bandūq thī.}
\end{quote}

\begin{quote}
One of the bandits had a gun in his hand.
\end{quote}

Feminine.

\begin{cousinweb}
\textbf{\hi{बंदूक़}} (\emph{bandūq}) is usually traced, through Persian, to Arabic \emph{bunduq}, a hazelnut and then a small round pellet; the step from pellet to gun is said to have come through the pellet-bow. Hindi also writes it without the dot.
\end{cousinweb}

\begin{grammarlens}[title={What you've built}]
One more word for telling what happened.
\end{grammarlens}

\subsection*{Your turn}
\begin{itemize}
\raggedright
  \item \emph{Say it:} \emph{bandūq}
  \item \emph{Say it:} \emph{bandūq}, once more
  \item \emph{Say it:} \emph{ek ḍākū ke hāth meṁ bandūq thī.} — One of the bandits had a gun in his hand.
  \item \emph{Recall:} say the Hindi for to rob, then the Hindi for a robbery (by a gang), then say \emph{bandūq} again
\end{itemize}

\begin{tcolorbox}[breakable,colback=teal!4,colframe=teal!35!black,title={Before you move on}]
What does \hi{बंदूक़} mean? (\textquotedblleft{}A gun\textquotedblright{}.) Say it once more.
\end{tcolorbox}
\section[\texorpdfstring{\hi{बाँधना} (bā̃dhnā)}{बाँधना (bā̃dhnā)}]{\hi{बाँधना} (bā̃dhnā) — to tie up; to bind}

\label{lesson:HI-C292-badhna}



\begin{quote}
Before the new one: say the Hindi for a robbery (by a gang), then the Hindi for a gun.
\end{quote}

\subsection*{You'll want to know: \hi{बाँधना}}
\textbf{\hi{बाँधना}} — \emph{bā̃dhnā} — \textquotedblleft{}to tie up; to bind\textquotedblright{}.

\begin{quote}
\hi{डाकुओं} \hi{ने} \hi{चौकीदार} \hi{को} \hi{रस्सी} \hi{से} \hi{बाँधा।}
\end{quote}

\begin{quote}
\emph{ḍākuoṁ ne caukīdār ko rassī se bā̃dhā.}
\end{quote}

\begin{quote}
The bandits tied the watchman up with rope.
\end{quote}

Past, with \hi{ने} (he or she …): \hi{उसने} \hi{रस्सी} \hi{बाँधी}, \emph{usne rassī bā̃dhī}.

\begin{cousinweb}
\textbf{\hi{बाँधना}} (\emph{bā̃dhnā}) goes back to Sanskrit \emph{bandhati}, \textquotedblleft{}binds\textquotedblright{} (CDIAL 9139). A person tied up takes \textbf{\hi{को}} and the verb stays masculine: \textbf{\hi{चौकीदार} \hi{को} \hi{बाँधा}}; a thing tied agrees: \textbf{\hi{रस्सी} \hi{बाँधी}}.
\end{cousinweb}

\begin{grammarlens}[title={What you've built}]
One more word for telling what happened.
\end{grammarlens}

\subsection*{Your turn}
\begin{itemize}
\raggedright
  \item \emph{Say it:} \emph{bā̃dhnā}
  \item \emph{Say it:} \emph{bā̃dhnā}, once more
  \item \emph{Say it:} \emph{ḍākuoṁ ne caukīdār ko rassī se bā̃dhā.} — The bandits tied the watchman up with rope.
  \item \emph{Recall:} say the Hindi for a robbery (by a gang), then the Hindi for a gun, then say \emph{bā̃dhnā} again
\end{itemize}

\begin{tcolorbox}[breakable,colback=teal!4,colframe=teal!35!black,title={Before you move on}]
What does \hi{बाँधना} mean? (\textquotedblleft{}To tie up; to bind\textquotedblright{}.) Say it once more.
\end{tcolorbox}
